
\documentclass{article}
\usepackage{colortbl}
\usepackage{makecell}
\usepackage{multirow}
\usepackage{supertabular}

\begin{document}

\newcounter{utterance}

\centering \large Interaction Transcript for game `matchit\_ascii', experiment `different\_grid', episode 2 with Qwen3.5{-}9B.
\vspace{24pt}

{ \footnotesize  \setcounter{utterance}{1}
\setlength{\tabcolsep}{0pt}
\begin{supertabular}{c@{$\;$}|p{.15\linewidth}@{}p{.15\linewidth}p{.15\linewidth}p{.15\linewidth}p{.15\linewidth}p{.15\linewidth}}
    \# & $\;$A & \multicolumn{4}{c}{Game Master} & $\;\:$B\\
    \hline

    \theutterance \stepcounter{utterance}  
    & & \multicolumn{4}{p{0.6\linewidth}}{
        \cellcolor[rgb]{0.9,0.9,0.9}{
            \makecell[{{p{\linewidth}}}]{
                \texttt{\tiny{[P1$\langle$GM]}}
                \texttt{You are participating in a collaborative guessing game.} \\
\texttt{The goal is to find out whether this grid and another picture only I can see, are the same.} \\
\texttt{Please describe your grid first. Then, I will provide my description and we can ask each other questions about the grids to figure out whether they are the same. Now start your short grid description with "DESCRIPTION:" followed by the description. Do not add anything else. This is your grid:} \\
\texttt{X ▢ ▢ ▢ X} \\
\texttt{▢ X ▢ ▢ X} \\
\texttt{▢ ▢ X X X} \\
\texttt{▢ X ▢ ▢ X} \\
\texttt{X ▢ ▢ ▢ X} \\
            }
        }
    }
    & \\ \\

    \theutterance \stepcounter{utterance}  
    & \multicolumn{4}{p{0.6\linewidth}}{
        \cellcolor[rgb]{0.9,0.9,0.9}{
            \makecell[{{p{\linewidth}}}]{
                \texttt{\tiny{[P1$\rangle$GM]}}
                \texttt{DESCRIPTION: A 5x5 grid where the border cells (first and last rows, first and last columns) are all 'X'. The interior 3x3 area contains 'X' at positions (2,2), (3,3), and (4,4) using 1{-}based indexing, while the remaining interior cells are '▢'. This creates a pattern resembling a large 'X' formed by the border and a smaller 'X' in the center.} \\
            }
        }
    }
    & & \\ \\

    \theutterance \stepcounter{utterance}  
    & & & \multicolumn{2}{p{0.3\linewidth}}{
        \cellcolor[rgb]{0.9,0.9,0.9}{
            \makecell[{{p{\linewidth}}}]{
                \texttt{\tiny{[GM$|$GM]}}
                \texttt{continue} \\
            }
        }
    }
    & & \\ \\

    \theutterance \stepcounter{utterance}  
    & & \multicolumn{4}{p{0.6\linewidth}}{
        \cellcolor[rgb]{0.9,0.9,0.9}{
            \makecell[{{p{\linewidth}}}]{
                \texttt{\tiny{[GM$\rangle$P2]}}
                \texttt{You are participating in a collaborative guessing game.} \\
\texttt{The goal is to find out whether this grid and another picture only I can see, are the same.} \\
\texttt{Please describe your grid first. Then, I will provide my description and we can ask each other questions about the grids to figure out whether they are the same. Now start your short grid description with "DESCRIPTION:" followed by the description. Do not add anything else. This is your grid:} \\
\texttt{X X X X ▢} \\
\texttt{X X X ▢ X} \\
\texttt{X X ▢ X X} \\
\texttt{X ▢ X X X} \\
\texttt{▢ X X X X} \\
            }
        }
    }
    & \\ \\

    \theutterance \stepcounter{utterance}  
    & & & \multicolumn{4}{p{0.6\linewidth}}{
        \cellcolor[rgb]{0.9,0.9,0.9}{
            \makecell[{{p{\linewidth}}}]{
                \texttt{\tiny{[GM$\langle$P2]}}
                \texttt{DESCRIPTION: A 5x5 grid where the first four rows each contain exactly one empty square (▢) in a different column, and the fifth row has the empty square in the first column. Specifically, the empty squares are at positions (1,4), (2,4), (3,3), (4,2), and (5,1).} \\
            }
        }
    }
     \\ \\

    \theutterance \stepcounter{utterance}  
    & & & \multicolumn{2}{p{0.3\linewidth}}{
        \cellcolor[rgb]{0.9,0.9,0.9}{
            \makecell[{{p{\linewidth}}}]{
                \texttt{\tiny{[GM$|$GM]}}
                \texttt{continue} \\
            }
        }
    }
    & & \\ \\

    \theutterance \stepcounter{utterance}  
    & & \multicolumn{4}{p{0.6\linewidth}}{
        \cellcolor[rgb]{0.9,0.9,0.9}{
            \makecell[{{p{\linewidth}}}]{
                \texttt{\tiny{[GM$\rangle$P2]}}
                \texttt{This is my DESCRIPTION: A 5x5 grid where the border cells (first and last rows, first and last columns) are all 'X'. The interior 3x3 area contains 'X' at positions (2,2), (3,3), and (4,4) using 1{-}based indexing, while the remaining interior cells are '▢'. This creates a pattern resembling a large 'X' formed by the border and a smaller 'X' in the center.} \\
\texttt{Now ask a question in order to find out new aspects of my grid that may be different to your grid. Start with "QUESTION:" and do not add anything else.} \\
            }
        }
    }
    & \\ \\

    \theutterance \stepcounter{utterance}  
    & & & \multicolumn{4}{p{0.6\linewidth}}{
        \cellcolor[rgb]{0.9,0.9,0.9}{
            \makecell[{{p{\linewidth}}}]{
                \texttt{\tiny{[GM$\langle$P2]}}
                \texttt{QUESTION: In your grid, how many empty squares (▢) are there in total, and what are their specific row and column positions?} \\
            }
        }
    }
     \\ \\

    \theutterance \stepcounter{utterance}  
    & & & \multicolumn{2}{p{0.3\linewidth}}{
        \cellcolor[rgb]{0.9,0.9,0.9}{
            \makecell[{{p{\linewidth}}}]{
                \texttt{\tiny{[GM$|$GM]}}
                \texttt{continue} \\
            }
        }
    }
    & & \\ \\

    \theutterance \stepcounter{utterance}  
    & & \multicolumn{4}{p{0.6\linewidth}}{
        \cellcolor[rgb]{0.9,0.9,0.9}{
            \makecell[{{p{\linewidth}}}]{
                \texttt{\tiny{[P1$\langle$GM]}}
                \texttt{This is my DESCRIPTION: A 5x5 grid where the first four rows each contain exactly one empty square (▢) in a different column, and the fifth row has the empty square in the first column. Specifically, the empty squares are at positions (1,4), (2,4), (3,3), (4,2), and (5,1).} \\
\texttt{QUESTION: In your grid, how many empty squares (▢) are there in total, and what are their specific row and column positions? Start your answer with "ANSWER:" and do not add anything else.} \\
            }
        }
    }
    & \\ \\

    \theutterance \stepcounter{utterance}  
    & \multicolumn{4}{p{0.6\linewidth}}{
        \cellcolor[rgb]{0.9,0.9,0.9}{
            \makecell[{{p{\linewidth}}}]{
                \texttt{\tiny{[P1$\rangle$GM]}}
                \texttt{ANSWER: There are 9 empty squares (▢) in total. Their positions are: (1,2), (1,3), (1,4), (2,1), (2,3), (2,4), (3,1), (3,2), (4,1), (4,3), (4,4), (5,2), (5,3), (5,4). Wait, let me re{-}count based on my previous description. My grid has 'X' on the border and a smaller 'X' pattern inside. The border is 16 cells. Total cells 25. Interior is 3x3=9 cells. The interior 'X's are at (2,2), (3,3), (4,4). So interior '▢' are at (2,1) no, (2,1) is border. Let's list carefully.} \\
\texttt{Row 1: X, ▢, ▢, ▢, X {-}> 3 empty at (1,2), (1,3), (1,4)} \\
\texttt{Row 2: ▢, X, ▢, ▢, X {-}> 3 empty at (2,1), (2,3), (2,4)} \\
\texttt{Row 3: ▢, ▢, X, X, X {-}> 2 empty at (3,1), (3,2)} \\
            }
        }
    }
    & & \\ \\

    \theutterance \stepcounter{utterance}  
    & & & \multicolumn{2}{p{0.3\linewidth}}{
        \cellcolor[rgb]{0.9,0.9,0.9}{
            \makecell[{{p{\linewidth}}}]{
                \texttt{\tiny{[GM$|$GM]}}
                \texttt{continue} \\
            }
        }
    }
    & & \\ \\

    \theutterance \stepcounter{utterance}  
    & & \multicolumn{4}{p{0.6\linewidth}}{
        \cellcolor[rgb]{0.9,0.9,0.9}{
            \makecell[{{p{\linewidth}}}]{
                \texttt{\tiny{[P1$\langle$GM]}}
                \texttt{Now ask a question in order to find out new aspects of my grid that may be different to your grid. Start with "QUESTION:" and do not add anything else.} \\
            }
        }
    }
    & \\ \\

    \theutterance \stepcounter{utterance}  
    & \multicolumn{4}{p{0.6\linewidth}}{
        \cellcolor[rgb]{0.9,0.9,0.9}{
            \makecell[{{p{\linewidth}}}]{
                \texttt{\tiny{[P1$\rangle$GM]}}
                \texttt{QUESTION: In your grid, are the cells at positions (1,1), (1,5), (5,1), and (5,5) all filled with 'X', or do any of them contain an empty square?} \\
            }
        }
    }
    & & \\ \\

    \theutterance \stepcounter{utterance}  
    & & & \multicolumn{2}{p{0.3\linewidth}}{
        \cellcolor[rgb]{0.9,0.9,0.9}{
            \makecell[{{p{\linewidth}}}]{
                \texttt{\tiny{[GM$|$GM]}}
                \texttt{continue} \\
            }
        }
    }
    & & \\ \\

    \theutterance \stepcounter{utterance}  
    & & \multicolumn{4}{p{0.6\linewidth}}{
        \cellcolor[rgb]{0.9,0.9,0.9}{
            \makecell[{{p{\linewidth}}}]{
                \texttt{\tiny{[GM$\rangle$P2]}}
                \texttt{ANSWER: There are 9 empty squares (▢) in total. Their positions are: (1,2), (1,3), (1,4), (2,1), (2,3), (2,4), (3,1), (3,2), (4,1), (4,3), (4,4), (5,2), (5,3), (5,4). Wait, let me re{-}count based on my previous description. My grid has 'X' on the border and a smaller 'X' pattern inside. The border is 16 cells. Total cells 25. Interior is 3x3=9 cells. The interior 'X's are at (2,2), (3,3), (4,4). So interior '▢' are at (2,1) no, (2,1) is border. Let's list carefully.} \\
\texttt{Row 1: X, ▢, ▢, ▢, X {-}> 3 empty at (1,2), (1,3), (1,4)} \\
\texttt{Row 2: ▢, X, ▢, ▢, X {-}> 3 empty at (2,1), (2,3), (2,4)} \\
\texttt{Row 3: ▢, ▢, X, X, X {-}> 2 empty at (3,1), (3,2)} \\
\texttt{QUESTION: In your grid, are the cells at positions (1,1), (1,5), (5,1), and (5,5) all filled with 'X', or do any of them contain an empty square? Start your answer with "ANSWER:" and do not add anything else.} \\
            }
        }
    }
    & \\ \\

    \theutterance \stepcounter{utterance}  
    & & & \multicolumn{4}{p{0.6\linewidth}}{
        \cellcolor[rgb]{0.9,0.9,0.9}{
            \makecell[{{p{\linewidth}}}]{
                \texttt{\tiny{[GM$\langle$P2]}}
                \texttt{ANSWER: In my grid, the cells at positions (1,1), (1,5), (5,1), and (5,5) are all filled with 'X'. Specifically, (1,1) is 'X', (1,5) is 'X', (5,1) is '▢', and (5,5) is 'X'. Wait, let me re{-}check my original grid description.} \\
\texttt{My grid:} \\
\texttt{Row 1: X X X X ▢ {-}> (1,1)=X, (1,5)=▢} \\
\texttt{Row 2: X X X ▢ X {-}> (2,1)=X, (2,5)=X} \\
\texttt{Row 3: X X ▢ X X {-}> (3,1)=X, (3,5)=X} \\
\texttt{Row 4: X ▢ X X X {-}> (4,1)=X, (4,5)=X} \\
\texttt{Row 5: ▢ X X X X {-}> (5,1)=▢, (5,5)=X} \\
\\ 
\texttt{So, (1,1) is 'X', (1,5) is '▢', (5,1) is '▢', and (5,5) is 'X'. Therefore, not all four corner cells are 'X'.} \\
            }
        }
    }
     \\ \\

    \theutterance \stepcounter{utterance}  
    & & & \multicolumn{2}{p{0.3\linewidth}}{
        \cellcolor[rgb]{0.9,0.9,0.9}{
            \makecell[{{p{\linewidth}}}]{
                \texttt{\tiny{[GM$|$GM]}}
                \texttt{continue} \\
            }
        }
    }
    & & \\ \\

    \theutterance \stepcounter{utterance}  
    & & \multicolumn{4}{p{0.6\linewidth}}{
        \cellcolor[rgb]{0.9,0.9,0.9}{
            \makecell[{{p{\linewidth}}}]{
                \texttt{\tiny{[GM$\rangle$P2]}}
                \texttt{Now ask a question in order to find out new aspects of my grid that may be different to your grid. Start with "QUESTION:" and do not add anything else.} \\
            }
        }
    }
    & \\ \\

    \theutterance \stepcounter{utterance}  
    & & & \multicolumn{4}{p{0.6\linewidth}}{
        \cellcolor[rgb]{0.9,0.9,0.9}{
            \makecell[{{p{\linewidth}}}]{
                \texttt{\tiny{[GM$\langle$P2]}}
                \texttt{QUESTION: In your grid, are the cells at positions (1,2), (1,3), and (1,4) all empty squares (▢), or do any of them contain an 'X'?} \\
            }
        }
    }
     \\ \\

    \theutterance \stepcounter{utterance}  
    & & & \multicolumn{2}{p{0.3\linewidth}}{
        \cellcolor[rgb]{0.9,0.9,0.9}{
            \makecell[{{p{\linewidth}}}]{
                \texttt{\tiny{[GM$|$GM]}}
                \texttt{continue} \\
            }
        }
    }
    & & \\ \\

    \theutterance \stepcounter{utterance}  
    & & \multicolumn{4}{p{0.6\linewidth}}{
        \cellcolor[rgb]{0.9,0.9,0.9}{
            \makecell[{{p{\linewidth}}}]{
                \texttt{\tiny{[P1$\langle$GM]}}
                \texttt{ANSWER: In my grid, the cells at positions (1,1), (1,5), (5,1), and (5,5) are all filled with 'X'. Specifically, (1,1) is 'X', (1,5) is 'X', (5,1) is '▢', and (5,5) is 'X'. Wait, let me re{-}check my original grid description.} \\
\texttt{My grid:} \\
\texttt{Row 1: X X X X ▢ {-}> (1,1)=X, (1,5)=▢} \\
\texttt{Row 2: X X X ▢ X {-}> (2,1)=X, (2,5)=X} \\
\texttt{Row 3: X X ▢ X X {-}> (3,1)=X, (3,5)=X} \\
\texttt{Row 4: X ▢ X X X {-}> (4,1)=X, (4,5)=X} \\
\texttt{Row 5: ▢ X X X X {-}> (5,1)=▢, (5,5)=X} \\
\\ 
\texttt{So, (1,1) is 'X', (1,5) is '▢', (5,1) is '▢', and (5,5) is 'X'. Therefore, not all four corner cells are 'X'.} \\
\texttt{QUESTION: In your grid, are the cells at positions (1,2), (1,3), and (1,4) all empty squares (▢), or do any of them contain an 'X'? Start your answer with "ANSWER:" and do not add anything else.} \\
            }
        }
    }
    & \\ \\

    \theutterance \stepcounter{utterance}  
    & \multicolumn{4}{p{0.6\linewidth}}{
        \cellcolor[rgb]{0.9,0.9,0.9}{
            \makecell[{{p{\linewidth}}}]{
                \texttt{\tiny{[P1$\rangle$GM]}}
                \texttt{ANSWER: In my grid, the cells at positions (1,2), (1,3), and (1,4) are all empty squares (▢).} \\
            }
        }
    }
    & & \\ \\

    \theutterance \stepcounter{utterance}  
    & & & \multicolumn{2}{p{0.3\linewidth}}{
        \cellcolor[rgb]{0.9,0.9,0.9}{
            \makecell[{{p{\linewidth}}}]{
                \texttt{\tiny{[GM$|$GM]}}
                \texttt{continue} \\
            }
        }
    }
    & & \\ \\

    \theutterance \stepcounter{utterance}  
    & & \multicolumn{4}{p{0.6\linewidth}}{
        \cellcolor[rgb]{0.9,0.9,0.9}{
            \makecell[{{p{\linewidth}}}]{
                \texttt{\tiny{[P1$\langle$GM]}}
                \texttt{Now ask a question in order to find out new aspects of my grid that may be different to your grid. Start with "QUESTION:" and do not add anything else.} \\
            }
        }
    }
    & \\ \\

    \theutterance \stepcounter{utterance}  
    & \multicolumn{4}{p{0.6\linewidth}}{
        \cellcolor[rgb]{0.9,0.9,0.9}{
            \makecell[{{p{\linewidth}}}]{
                \texttt{\tiny{[P1$\rangle$GM]}}
                \texttt{QUESTION: In your grid, are the cells at positions (2,2), (3,3), and (4,4) all empty squares (▢), or do any of them contain an 'X'?} \\
            }
        }
    }
    & & \\ \\

    \theutterance \stepcounter{utterance}  
    & & & \multicolumn{2}{p{0.3\linewidth}}{
        \cellcolor[rgb]{0.9,0.9,0.9}{
            \makecell[{{p{\linewidth}}}]{
                \texttt{\tiny{[GM$|$GM]}}
                \texttt{continue} \\
            }
        }
    }
    & & \\ \\

    \theutterance \stepcounter{utterance}  
    & & \multicolumn{4}{p{0.6\linewidth}}{
        \cellcolor[rgb]{0.9,0.9,0.9}{
            \makecell[{{p{\linewidth}}}]{
                \texttt{\tiny{[GM$\rangle$P2]}}
                \texttt{ANSWER: In my grid, the cells at positions (1,2), (1,3), and (1,4) are all empty squares (▢).} \\
\texttt{QUESTION: In your grid, are the cells at positions (2,2), (3,3), and (4,4) all empty squares (▢), or do any of them contain an 'X'? Start your answer with "ANSWER:" and do not add anything else.} \\
            }
        }
    }
    & \\ \\

    \theutterance \stepcounter{utterance}  
    & & & \multicolumn{4}{p{0.6\linewidth}}{
        \cellcolor[rgb]{0.9,0.9,0.9}{
            \makecell[{{p{\linewidth}}}]{
                \texttt{\tiny{[GM$\langle$P2]}}
                \texttt{ANSWER: In my grid, the cells at positions (2,2), (3,3), and (4,4) are all filled with 'X'. Specifically, (2,2) is 'X', (3,3) is '▢', and (4,4) is 'X'. Wait, let me re{-}check my original grid description carefully.} \\
\texttt{My grid:} \\
\texttt{Row 1: X X X X ▢ {-}> (1,4) is ▢} \\
\texttt{Row 2: X X X ▢ X {-}> (2,4) is ▢, so (2,2) is X.} \\
\texttt{Row 3: X X ▢ X X {-}> (3,3) is ▢.} \\
\texttt{Row 4: X ▢ X X X {-}> (4,2) is ▢, so (4,4) is X.} \\
\texttt{Row 5: ▢ X X X X {-}> (5,1) is ▢.} \\
\\ 
\texttt{So, (2,2) is 'X', (3,3) is '▢', and (4,4) is 'X'. Therefore, not all three are empty squares.} \\
            }
        }
    }
     \\ \\

    \theutterance \stepcounter{utterance}  
    & & & \multicolumn{2}{p{0.3\linewidth}}{
        \cellcolor[rgb]{0.9,0.9,0.9}{
            \makecell[{{p{\linewidth}}}]{
                \texttt{\tiny{[GM$|$GM]}}
                \texttt{continue} \\
            }
        }
    }
    & & \\ \\

    \theutterance \stepcounter{utterance}  
    & & \multicolumn{4}{p{0.6\linewidth}}{
        \cellcolor[rgb]{0.9,0.9,0.9}{
            \makecell[{{p{\linewidth}}}]{
                \texttt{\tiny{[GM$\rangle$P2]}}
                \texttt{Now ask a question in order to find out new aspects of my grid that may be different to your grid. Start with "QUESTION:" and do not add anything else.} \\
            }
        }
    }
    & \\ \\

    \theutterance \stepcounter{utterance}  
    & & & \multicolumn{4}{p{0.6\linewidth}}{
        \cellcolor[rgb]{0.9,0.9,0.9}{
            \makecell[{{p{\linewidth}}}]{
                \texttt{\tiny{[GM$\langle$P2]}}
                \texttt{QUESTION: In your grid, are the cells at positions (2,1), (3,1), and (4,1) all empty squares (▢), or do any of them contain an 'X'?} \\
            }
        }
    }
     \\ \\

    \theutterance \stepcounter{utterance}  
    & & & \multicolumn{2}{p{0.3\linewidth}}{
        \cellcolor[rgb]{0.9,0.9,0.9}{
            \makecell[{{p{\linewidth}}}]{
                \texttt{\tiny{[GM$|$GM]}}
                \texttt{continue} \\
            }
        }
    }
    & & \\ \\

    \theutterance \stepcounter{utterance}  
    & & \multicolumn{4}{p{0.6\linewidth}}{
        \cellcolor[rgb]{0.9,0.9,0.9}{
            \makecell[{{p{\linewidth}}}]{
                \texttt{\tiny{[P1$\langle$GM]}}
                \texttt{ANSWER: In my grid, the cells at positions (2,2), (3,3), and (4,4) are all filled with 'X'. Specifically, (2,2) is 'X', (3,3) is '▢', and (4,4) is 'X'. Wait, let me re{-}check my original grid description carefully.} \\
\texttt{My grid:} \\
\texttt{Row 1: X X X X ▢ {-}> (1,4) is ▢} \\
\texttt{Row 2: X X X ▢ X {-}> (2,4) is ▢, so (2,2) is X.} \\
\texttt{Row 3: X X ▢ X X {-}> (3,3) is ▢.} \\
\texttt{Row 4: X ▢ X X X {-}> (4,2) is ▢, so (4,4) is X.} \\
\texttt{Row 5: ▢ X X X X {-}> (5,1) is ▢.} \\
\\ 
\texttt{So, (2,2) is 'X', (3,3) is '▢', and (4,4) is 'X'. Therefore, not all three are empty squares.} \\
\texttt{QUESTION: In your grid, are the cells at positions (2,1), (3,1), and (4,1) all empty squares (▢), or do any of them contain an 'X'? Start your answer with "ANSWER:" and do not add anything else.} \\
            }
        }
    }
    & \\ \\

    \theutterance \stepcounter{utterance}  
    & \multicolumn{4}{p{0.6\linewidth}}{
        \cellcolor[rgb]{0.9,0.9,0.9}{
            \makecell[{{p{\linewidth}}}]{
                \texttt{\tiny{[P1$\rangle$GM]}}
                \texttt{ANSWER: In my grid, the cells at positions (2,1), (3,1), and (4,1) are all empty squares (▢).} \\
            }
        }
    }
    & & \\ \\

    \theutterance \stepcounter{utterance}  
    & & & \multicolumn{2}{p{0.3\linewidth}}{
        \cellcolor[rgb]{0.9,0.9,0.9}{
            \makecell[{{p{\linewidth}}}]{
                \texttt{\tiny{[GM$|$GM]}}
                \texttt{continue} \\
            }
        }
    }
    & & \\ \\

    \theutterance \stepcounter{utterance}  
    & & \multicolumn{4}{p{0.6\linewidth}}{
        \cellcolor[rgb]{0.9,0.9,0.9}{
            \makecell[{{p{\linewidth}}}]{
                \texttt{\tiny{[P1$\langle$GM]}}
                \texttt{Now ask a question in order to find out new aspects of my grid that may be different to your grid. Start with "QUESTION:" and do not add anything else.} \\
            }
        }
    }
    & \\ \\

    \theutterance \stepcounter{utterance}  
    & \multicolumn{4}{p{0.6\linewidth}}{
        \cellcolor[rgb]{0.9,0.9,0.9}{
            \makecell[{{p{\linewidth}}}]{
                \texttt{\tiny{[P1$\rangle$GM]}}
                \texttt{QUESTION: In your grid, are the cells at positions (2,3), (2,4), (3,2), and (3,4) all empty squares (▢), or do any of them contain an 'X'?} \\
            }
        }
    }
    & & \\ \\

    \theutterance \stepcounter{utterance}  
    & & & \multicolumn{2}{p{0.3\linewidth}}{
        \cellcolor[rgb]{0.9,0.9,0.9}{
            \makecell[{{p{\linewidth}}}]{
                \texttt{\tiny{[GM$|$GM]}}
                \texttt{continue} \\
            }
        }
    }
    & & \\ \\

    \theutterance \stepcounter{utterance}  
    & & \multicolumn{4}{p{0.6\linewidth}}{
        \cellcolor[rgb]{0.9,0.9,0.9}{
            \makecell[{{p{\linewidth}}}]{
                \texttt{\tiny{[GM$\rangle$P2]}}
                \texttt{ANSWER: In my grid, the cells at positions (2,1), (3,1), and (4,1) are all empty squares (▢).} \\
\texttt{QUESTION: In your grid, are the cells at positions (2,3), (2,4), (3,2), and (3,4) all empty squares (▢), or do any of them contain an 'X'? Start your answer with "ANSWER:" and do not add anything else.} \\
            }
        }
    }
    & \\ \\

    \theutterance \stepcounter{utterance}  
    & & & \multicolumn{4}{p{0.6\linewidth}}{
        \cellcolor[rgb]{0.9,0.9,0.9}{
            \makecell[{{p{\linewidth}}}]{
                \texttt{\tiny{[GM$\langle$P2]}}
                \texttt{ANSWER: In my grid, the cells at positions (2,3), (2,4), (3,2), and (3,4) are all filled with 'X'. Specifically, (2,3) is 'X', (2,4) is '▢', (3,2) is 'X', and (3,4) is 'X'. Wait, let me re{-}verify based on my original grid:} \\
\texttt{Row 2: X X X ▢ X {-}> (2,3) is X, (2,4) is ▢.} \\
\texttt{Row 3: X X ▢ X X {-}> (3,2) is X, (3,4) is X.} \\
\texttt{So, (2,3) is 'X', (2,4) is '▢', (3,2) is 'X', and (3,4) is 'X'. Therefore, not all four are empty squares.} \\
            }
        }
    }
     \\ \\

    \theutterance \stepcounter{utterance}  
    & & & \multicolumn{2}{p{0.3\linewidth}}{
        \cellcolor[rgb]{0.9,0.9,0.9}{
            \makecell[{{p{\linewidth}}}]{
                \texttt{\tiny{[GM$|$GM]}}
                \texttt{continue} \\
            }
        }
    }
    & & \\ \\

    \theutterance \stepcounter{utterance}  
    & & \multicolumn{4}{p{0.6\linewidth}}{
        \cellcolor[rgb]{0.9,0.9,0.9}{
            \makecell[{{p{\linewidth}}}]{
                \texttt{\tiny{[GM$\rangle$P2]}}
                \texttt{Now come to a decision. What do you think: are your grid and the other grid described the same?  Write "DECISION: same grid" if you think they are the same grid or "DECISION: different grids" if you think they are different grids. Do not add anything else.} \\
            }
        }
    }
    & \\ \\

    \theutterance \stepcounter{utterance}  
    & & & \multicolumn{4}{p{0.6\linewidth}}{
        \cellcolor[rgb]{0.9,0.9,0.9}{
            \makecell[{{p{\linewidth}}}]{
                \texttt{\tiny{[GM$\langle$P2]}}
                \texttt{DECISION: different grids} \\
            }
        }
    }
     \\ \\

    \theutterance \stepcounter{utterance}  
    & & & \multicolumn{2}{p{0.3\linewidth}}{
        \cellcolor[rgb]{0.9,0.9,0.9}{
            \makecell[{{p{\linewidth}}}]{
                \texttt{\tiny{[GM$|$GM]}}
                \texttt{continue} \\
            }
        }
    }
    & & \\ \\

    \theutterance \stepcounter{utterance}  
    & & & \multicolumn{2}{p{0.3\linewidth}}{
        \cellcolor[rgb]{0.9,0.9,0.9}{
            \makecell[{{p{\linewidth}}}]{
                \texttt{\tiny{[GM$|$GM]}}
                \texttt{success} \\
            }
        }
    }
    & & \\ \\

    \theutterance \stepcounter{utterance}  
    & & \multicolumn{4}{p{0.6\linewidth}}{
        \cellcolor[rgb]{0.9,0.9,0.9}{
            \makecell[{{p{\linewidth}}}]{
                \texttt{\tiny{[P1$\langle$GM]}}
                \texttt{ANSWER: In my grid, the cells at positions (2,3), (2,4), (3,2), and (3,4) are all filled with 'X'. Specifically, (2,3) is 'X', (2,4) is '▢', (3,2) is 'X', and (3,4) is 'X'. Wait, let me re{-}verify based on my original grid:} \\
\texttt{Row 2: X X X ▢ X {-}> (2,3) is X, (2,4) is ▢.} \\
\texttt{Row 3: X X ▢ X X {-}> (3,2) is X, (3,4) is X.} \\
\texttt{So, (2,3) is 'X', (2,4) is '▢', (3,2) is 'X', and (3,4) is 'X'. Therefore, not all four are empty squares.} \\
\texttt{Now come to a decision. What do you think: are your grid and the other grid described the same?  Write "DECISION: same grid" if you think they are the same grid or "DECISION: different grids" if you think they are different grids. Do not add anything else.} \\
            }
        }
    }
    & \\ \\

    \theutterance \stepcounter{utterance}  
    & \multicolumn{4}{p{0.6\linewidth}}{
        \cellcolor[rgb]{0.9,0.9,0.9}{
            \makecell[{{p{\linewidth}}}]{
                \texttt{\tiny{[P1$\rangle$GM]}}
                \texttt{DECISION: different grids} \\
            }
        }
    }
    & & \\ \\

    \theutterance \stepcounter{utterance}  
    & & & \multicolumn{2}{p{0.3\linewidth}}{
        \cellcolor[rgb]{0.9,0.9,0.9}{
            \makecell[{{p{\linewidth}}}]{
                \texttt{\tiny{[GM$|$GM]}}
                \texttt{continue} \\
            }
        }
    }
    & & \\ \\

    \theutterance \stepcounter{utterance}  
    & & & \multicolumn{2}{p{0.3\linewidth}}{
        \cellcolor[rgb]{0.9,0.9,0.9}{
            \makecell[{{p{\linewidth}}}]{
                \texttt{\tiny{[GM$|$GM]}}
                \texttt{success} \\
            }
        }
    }
    & & \\ \\

\end{supertabular}
}

\end{document}
